\section{Refinement Experiments}
\label{sec:g-series}

The cascade as specified in Section~\ref{sec:tch} was the result of two rounds of experimentation. The first round (the \emph{baseline ablations}) was an exhaustive 13-experiment optimization that locked the cascade as ``v2f''. The second round (the \emph{G-series}) was eight targeted late-stage refinements that took us from v2f to the final locked headline. We document both because the negative results in each are as informative as the positives.

\subsection{Baseline ablations (v2f lock-in)}
\label{sec:g-series-ablations}

Thirteen independent ablations were run against the v2f cascade. Each tested one design dimension; each kept or rejected its variant based on the measured delta. The unified conclusion is that the cascade in its final form is at a measured local optimum of the metric surface around it.

\begin{table}[h]
\centering
\caption{13 baseline ablations on the v2f cascade. Outcomes that improved the cascade are \textsc{kept}; outcomes that hurt are \textsc{rejected} (the simpler design wins).}
\label{tab:baseline-ablations}
\small
\setlength{\tabcolsep}{4pt}
\renewcommand{\arraystretch}{1.15}
\begin{tabularx}{\linewidth}{@{}r >{\raggedright\arraybackslash}X >{\raggedright\arraybackslash}X@{}}
\toprule
\# & Question & Outcome \\
\midrule
1 & Which retrievers to include in L2 RRF? Drop-one over the six. & \textsc{Drop} \hrrf{} LLM ($+2.1$~pts Hit@5). Keep the other four. \\
2 & Score-weighted RRF (weight by standalone Hit@5)? & All variants hurt by 4--6 points Hit@5; \textsc{keep} uniform weights. \\
3 & Per-family oracle (knew gold family, pick best retriever)? & Oracle Hit@5 $= 0.888$, \emph{worse} than uniform RRF $0.912$; \textsc{reject}. \\
4 & Confidence-aware RRF (weight by per-window triage)? & Hurts by 5--7 points Hit@1; \textsc{reject}. \\
5 & Learned per-candidate ranker (LogReg on rank features)? & Hit@5 $= 0.825$, far below cascade $0.912$; \textsc{reject}. \\
6 & Overlap-rerank window sweep (top-3 vs top-5 vs top-10)? & top-3 + 3-voter pool optimal at Hit@1 $= 0.722$; \textsc{keep top-3}. \\
7 & Let \agent{} override L2 top-1? & Net $-5$ Hit@1 wins on 200 windows; \textsc{disable} agent rerank. \\
8 & L4 stacker: GBM vs LogReg? & LogReg PR-AUC $0.9998 / 0.853$ vs GBM $0.985 / 0.739$; \textsc{keep} LogReg. \\
9 & L4 derived features (max-conf, variance, consensus)? & Inclusive PR-AUC drops $0.853 \to 0.849$; \textsc{reject}. \\
10 & L4 without \hgb{}? & Strict PR-AUC collapses $0.9998 \to 0.303$; \hgb{} is essential. \\
11 & L1 noise threshold sweep $\{0.1, 0.2, 0.3, 0.5\}$? & Recall stays $1.000$, F1 peaks at $\tau = 0.5$; \textsc{use} 0.5. \\
12 & Free novelty signal vs agent's \texttt{is\_novel}? & At $\tau = 0.5$: 96\% precision (matches agent); OR-combine widens recall; \textsc{use combined OR}. \\
13 & Multi-seed stability of stacker (seeds 1, 7, 42, 100, 1000)? & Strict PR-AUC std $= 0.0000$; inclusive std $= 0.0025$; \textsc{not} a fortunate seed. \\
\bottomrule
\end{tabularx}
\end{table}

The 13 ablations confirmed the local optimum. We also computed the \emph{theoretical ceiling}: the union of the four L2 retrievers' top-5 sets contains the gold ticket on $97.6\%$ of windows---the absolute maximum any fusion of these specific retrievers can achieve. The v2f cascade at Hit@5 $= 0.912$ is at $93.5\%$ of that ceiling. Closing the remaining 6.4-point gap requires a \emph{new} retriever, not further fusion tuning. This conclusion is what motivated the G-series.

\subsection{The G-series: eight targeted refinements}

The G-series was launched on 2026-06-05 with a single explicit purpose: try to close part of the 6.4-point Hit@5 gap to the union ceiling, or, failing that, find improvements on the other axes (novelty recall, robustness, OOD generalization). Phases were sequential; each was committed independently with its own observation document under \texttt{docs4/}. Table~\ref{tab:g-summary} summarizes outcomes.

\begin{table}[h]
\centering
\caption{G-series outcomes. KEEP-headline = integrated into the final locked cascade. KEEP-analysis = the experiment is published as a robustness/generalization result but does not modify the cascade. SKIP = the variant hurt or didn't help.}
\label{tab:g-summary}
\small
\setlength{\tabcolsep}{4pt}
\renewcommand{\arraystretch}{1.15}
\begin{tabularx}{\linewidth}{@{}l >{\raggedright\arraybackslash}X >{\raggedright\arraybackslash}X@{}}
\toprule
Phase & Intervention & Verdict \\
\midrule
G1 & BiEncoder fine-tune with BM25 hard-negs + random negs & \textbf{KEEP} (Hit@1 $+2.1\%$ rel) \\
G2 & Fine-tuned cross-encoder as 5th retriever / reranker & SKIP (Hit@1 $-5.4$~pts) \\
G3 & Symmetric LLM extraction (windows too, not just tickets) & SKIP (standalone Hit@1 $+404\%$ rel; cascade integration breaks) \\
G4 & \agent{} Phase 3: complete coverage on remaining 658 windows & \textbf{KEEP} (novel recall $+119\%$ rel) \\
G5 & LLM judge reranker (Qwen picks best of top-5) & SKIP (third reranker negative; uninformative confidence) \\
G6 & Distractor robustness sweep (simulated, $r = 0, 10, 25, 50\%$) & KEEP-analysis (Hit@5 $-2\%$ at $r=50\%$) \\
G7 & Learned per-window novelty threshold (LogReg, 46 features) & \textbf{KEEP} (novel recall $+388\%$ rel; the headline win) \\
G8 & OOD evaluation: leave-one-family-out cross-validation & KEEP-analysis (F1 within $-11.4\%$ rel of in-distribution) \\
\bottomrule
\end{tabularx}
\end{table}

\subsubsection{G1 --- BiEncoder fine-tune with mixed negatives (KEEP)}

\textbf{Hypothesis.} The locked v2 \bienc{} fine-tune used only BM25-mined hard negatives, possibly over-fitting to lexical overlap. Mixing one random negative per anchor (sampled from memory tickets not in gold and not in BM25 top-20) should force semantic discrimination beyond lexical.

\textbf{Setup.} \texttt{sentence-transformers/all-MiniLM-L6-v2}, \texttt{MultipleNegativesRankingLoss}, $n_{\text{hard-neg}} = 2$, $n_{\text{random-neg}} = 1$ per anchor. 16{,}338 training examples from train+val. Five epochs, batch 32, $\mathrm{lr} = 2 \times 10^{-5}$ with 10\% warmup. Twenty-minute wall on RTX 5060.

\textbf{Result.} Hit@1 lifted from 0.7069 (v2f) to 0.7221 ($+2.1\%$ rel) at preserved Hit@5 (0.9124) and a tiny novelty regression ($-0.9\%$ rel). Inclusive PR-AUC ticked up $+0.4\%$ rel.

\textbf{Decision.} Integrated. G1's predictions file (\texttt{predictions-bienc.jsonl}) replaces the upstream \bienc{} source in every subsequent cascade build via the \texttt{TCH\_OVERRIDE\_BIENC} env variable.

\subsubsection{G2 --- Cross-encoder reranker (SKIP)}

\textbf{Hypothesis.} A cross-encoder reads query + document jointly (unlike bi-encoders, which encode separately) and should catch interactions a bi-encoder structurally cannot. We retrained \texttt{cross-encoder/ms-marco-MiniLM-L-6-v2} on the v2 split with BCE loss, 3 epochs, mixed negatives matching G1.

\textbf{Result.} Standalone training was healthy (F1 = 0.940, AP = 0.968 on val). But when integrated into the cascade in either of two natural modes---(a) added as a 5th L2 RRF voter, or (b) used as a reranker over the L2 top-5---it \emph{hurt} Hit@1 by 5.4 and 4.5 points respectively. Hit@5 dropped 1.8 points in mode (a) and was tied in mode (b).

\textbf{Why it failed.} The cross-encoder's standalone Hit@1 was 0.571, well below \bienc{}'s 0.722. Mode (a) suffered from the same ``RRF density paradox'' we saw with \hrrf{}-LLM. Mode (b) reordered the same 5 candidates and produced worse orderings than the cascade's existing anchor + overlap rerank.

\textbf{Decision.} SKIP. The negative result is publishable as supporting evidence for the cascade's local-optimum claim.

\subsubsection{G3 --- Symmetric LLM extraction (SKIP, but informative)}

\textbf{Hypothesis.} The asymmetric LLM-graph setup (tickets extracted by LLM, windows extracted by rule) leaves the kg retriever underperforming because the two sides emit different entity vocabularies. LLM-extracting the 1{,}008 test windows under the same strict JSON schema as the tickets should close that gap.

\textbf{Setup.} Qwen 3.6 35B-A3B via LM Studio, \texttt{WINDOW\_EXTRACTION\_SCHEMA}, thinking=OFF, $\sim$7.9 sec/window. 133 minutes wall time. 100\% extraction success.

\textbf{Standalone result (massive).} \kgret{} Hit@1 went from 0.078 (rule-extracted baseline) to 0.396 ($+404\%$ rel). Hit@5 went from 0.555 to 0.668 ($+20\%$). \hrrf{}-LLM with symmetric extraction: Hit@5 from 0.668 to 0.780 ($+17\%$). These are among the largest single-component lifts of the project.

\textbf{Cascade result (negative).} Every cascade configuration that integrated the G3 predictions hurt Hit@5 by 5\%+. We diagnosed two mechanisms:

\begin{enumerate}
\item \emph{Overlap-rerank vote shifts.} L2 position 1 depends on counting how many retrievers vote each \bienc{} top-3 candidate into their own top-3. G3 changes one retriever's top-3 sharply; ${\sim}99\%$ of windows had their top-1 vote distribution shift.
\item \emph{L4 stacker score-scale mismatch.} G3's confidence distribution differs from the asymmetric baseline; the 5-fold-CV LogReg refits but the L3 free-novelty trigger (\texttt{max\_conf < 0.5}) fires on a different population.
\end{enumerate}

\textbf{Decision.} SKIP from the cascade. The G3 artifacts (window extractions + predictions) are preserved for future score-scale-normalization experiments. The result is the cleanest demonstration in the project that \emph{the cascade is composition-fragile, not component-fragile}: a +404\% rel standalone win on one component breaks the cascade because the cascade's internal RRF + stacker math assumes specific score distributions.

\subsubsection{G4 --- DiagnosisAgent Phase 3 (KEEP)}

\textbf{Hypothesis.} The agent had run on 350 of 1{,}008 test windows (Phase 1: 200 random; Phase 2: 150 hard-case). Linear extrapolation from agent calibration ($\sim$92\% novelty precision, $\sim$30--40\% recall on the windows it had seen) predicted that running on the remaining 658 windows would lift novel recall from 0.162 to $\sim$0.30--0.35 at minor precision cost.

\textbf{Setup.} Same agent as Phase E. Six hours wall time on Qwen 35B with thinking=ON, $\sim$33 sec/window. One JSON parse failure handled with a novel fallback.

\textbf{Result.} Novel recall $0.162 \to 0.356$ ($+119\%$ rel). Novel precision $0.940 \to 0.930$ ($-1.0\%$ rel). Hit@1, Hit@5, MRR, and all PR-AUCs unchanged---the agent's outputs feed only L3, so they do not perturb retrieval or triage.

\textbf{Decision.} INTEGRATE. The G4 file replaces the agent source via \texttt{TCH\_EXTRA\_AGENT\_FILES} in subsequent builds. This was the first phase that strictly improved every headline metric (every metric tied or up, novel recall doubled).

\subsubsection{G5 --- LLM judge reranker (SKIP)}

\textbf{Hypothesis.} After L2 produces the top-5, call Qwen with thinking=ON and a strict schema (\texttt{\{best\_idx: 0--4, confidence: 0--1, reasoning\}}), then confidence-gate-override L2's top-1.

\textbf{Result.} Confidence is \emph{uninformative}: 84\% of 1{,}007 successful calls reported confidence in $[0.80, 0.90]$, with another 13\% at $\geq 0.90$ and only 3\% below. Threshold sweep over $\{0.5, 0.6, 0.7, 0.8, 0.9, 1.0\}$ showed every override level hurt Hit@1 vs the cascade. The best case (only override at confidence $\geq 0.9$, applied 90 times): Hit@1 dropped $0.7221 \to 0.6949$ ($-2.7$~pts). The worst case (override all 613 suggestions): $-13$ points.

\textbf{Decision.} SKIP. G5 makes the third independent reranker (after G2 cross-encoder and Phase E agent rerank) to underperform the cascade's overlap-rerank. The pattern is publishable.

\subsubsection{G6 --- Distractor robustness sweep (KEEP analysis)}

\textbf{Hypothesis.} Hit@5 should degrade gracefully as memory noise increases. We simulated distractor injection at the prediction level rather than re-fitting BiEncoder + SPLADE + KG per distractor ratio.

\textbf{Setup.} For each top-5 prediction, with per-slot probability $p = N_{\text{dist}} / (N_{\text{dist}} + N_{\text{memory}})$, replace the slot with a synthetic distractor ID (distractors are never gold by construction, so always count as misses). Ratios $r \in \{0, 10, 25, 50\}\%$ map to $p \in \{0, 0.031, 0.072, 0.137\}$.

\textbf{Result.} At $r = 50\%$: Hit@5 drops only $0.912 \to 0.894$ ($-2.0\%$ rel). Hit@1 drops $0.722 \to 0.619$ ($-14.2\%$ rel)---more sensitive but still competitive with the best single retriever's baseline Hit@1. MRR drops $-8.2\%$ rel. Sensitivity ranks: Hit@1 $\gg$ MRR $>$ Hit@5.

\textbf{Decision.} KEEP as an analysis result; the cascade is not modified. This is the simulation; an authentic re-fit sweep is left as future work (estimated 4 hours GPU, was not on the critical path).

\subsubsection{G7 --- Learned per-window novelty threshold (KEEP, headline win)}

\textbf{Motivation.} The L3 free-novelty signal (\texttt{max\_conf < 0.5}) is a single fixed threshold for all windows. Different families have different baseline retrieval confidences (active-fault windows naturally retrieve more confidently than baseline-traffic windows), so a fixed threshold under-flags some and over-flags others.

\textbf{Setup.} Logistic regression with \texttt{class\_weight=balanced} over 46 features: \texttt{window\_type} one-hot (4 levels), \texttt{scenario\_family} one-hot (27 levels), \texttt{service\_name} one-hot ($\sim$12), \texttt{is\_hard\_case}, \texttt{max\_retrieval\_conf}, \texttt{triage\_score}. Trained via 5-fold \texttt{StratifiedKFold} (\texttt{random\_state}=42) on the 1{,}008-window evaluation set, producing out-of-fold $P(\text{novel})$ per window.

\textbf{Leakage check.} A first iteration included \texttt{n\_prior\_family\_tickets} (count of past memory tickets in the same family). This gave F1 $= 1.000$---a near-perfect prediction. Investigation showed this is a tautology in our dataset: novelty $=$ ``no matching past memory ticket'', and $n_{\text{prior}} = 0$ effectively means the same thing once \texttt{scenario\_family} ties gold to memory. We excluded this feature for the headline model.

\textbf{Result (leakage-free, 5-fold out-of-fold).}

\begin{table}[h]
\centering
\caption{G7 leakage-free threshold sweep, out-of-fold predictions on 1{,}008 windows.}
\label{tab:g7-sweep}
\small
\begin{tabular}{rrrrrrr}
\toprule
$\tau$ & Flagged & TP & FP & Precision & Recall & F1 \\
\midrule
0.30 & 665 & 597 & 68 & 0.898 & 0.882 & \textbf{0.890} \\
0.40 & 594 & 556 & 38 & 0.936 & 0.821 & 0.875 \\
0.50 & 542 & 525 & 17 & 0.969 & 0.776 & 0.861 \\
0.60 & 526 & 521 & 5  & 0.991 & 0.770 & 0.866 \\
0.70 & 516 & 516 & 0  & 1.000 & 0.762 & 0.865 \\
\bottomrule
\end{tabular}
\end{table}

\textbf{Cascade integration.} At threshold $\tau = 0.50$ (chosen to match the v2f baseline's novel precision exactly), the cascade achieves novel precision $= 0.9405$ (tied with v2f) and novel recall $= 0.7932$ ($+388\%$ rel over the v2f baseline of 0.1625). Hit@1, Hit@5, MRR, and PR-AUC are unchanged---L3 only affects the novelty channel.

\textbf{Interpretation.} The top coefficients are exactly what one would expect:
\texttt{window\_type=pre\_fault\_baseline} ($+2.90$, novel by construction);
\texttt{window\_type=recovery\_window} ($-2.28$, post-incident, gold linked);
\texttt{scenario\_family=ad-outage} ($+2.62$, sparsely covered);
\texttt{scenario\_family=productcatalog-outage} ($-1.99$, well-covered).
G7 is exploiting structural metadata signal that the original fixed-threshold L3 ignored. The intervention is honestly a \emph{finding about the original cascade's blind spot}, not deep ML.

\textbf{Decision.} INTEGRATE at $\tau = 0.50$ as the headline win.

\subsubsection{G8 --- Out-of-distribution evaluation (KEEP analysis)}

\textbf{Motivation.} G7's in-distribution F1 numbers assume the test windows come from families seen during training. The harder question: does G7 generalize to families it has never seen?

\textbf{Setup.} Leave-one-family-out (LOFO) cross-validation: 27 folds, each holding out one of the 27 scenario families. Train the G7 LogReg on the remaining ${\sim}21$ families; predict on the held-out family. One-hot \texttt{scenario\_family=X} is always 0 in the held-out fold (since X was never seen during fit), so the model must use family-independent signals (\texttt{window\_type}, \texttt{max\_conf}, \texttt{triage\_score}, \texttt{is\_hard\_case}).

\textbf{Aggregate result (Table~\ref{tab:g8-sweep}).} OOD F1 is within $-11.4\%$ rel of in-distribution at the best threshold ($\tau = 0.30$).

\begin{table}[h]
\centering
\caption{G8 OOD threshold sweep vs in-distribution G7 (rows in $\Delta$ column are relative changes; negative means OOD is worse).}
\label{tab:g8-sweep}
\small
\begin{tabular}{rrrrrrr}
\toprule
$\tau$ & OOD prec & OOD rec & OOD F1 & ID F1 & $\Delta$ F1 rel \\
\midrule
0.30 & 0.793 & 0.783 & \textbf{0.788} & 0.890 & $-11.4\%$ \\
0.40 & 0.872 & 0.705 & 0.779 & 0.875 & $-10.9\%$ \\
0.50 & 0.903 & 0.616 & 0.732 & 0.861 & $-15.0\%$ \\
0.60 & 0.993 & 0.598 & 0.747 & 0.866 & $-13.8\%$ \\
0.70 & 1.000 & 0.597 & 0.748 & 0.865 & $-13.6\%$ \\
\bottomrule
\end{tabular}
\end{table}

\textbf{Per-family pattern.} \emph{Precision} is robust: 19 of 27 families exceed $0.90$ precision at $\tau = 0.50$. \emph{Recall} splits sharply by family structure:

\begin{itemize}
\item Families dominated by \texttt{window\_type=pre\_fault\_baseline} or \texttt{recovery\_window} generalize perfectly (\texttt{baseline-normal}: 100\%; \texttt{checkout-restart}: 100\%; \texttt{dns-outage}: 100\%; \texttt{slow-leak-saturation}: 100\%).
\item Families dominated by \texttt{active\_fault} windows lose 30--50~pts recall when held out (\texttt{ad-outage}: 31\%; \texttt{frontend-traffic-pressure}: 29\%; \texttt{email-outage}: 39\%).
\end{itemize}

\textbf{Interpretation.} The classifier learned a mostly-generic predictor with a family-specific boost. Hold the family out, lose the boost, and the generic component carries $\sim$80\% of the way. This is the correct OOD behavior; the model is not memorizing family identities.

\textbf{vs the v2f baseline at OOD.} Even at $\tau = 0.30$ OOD, the cascade lifts novel recall from the v2f baseline's $0.163$ to $0.783$ ($+381\%$ rel). The G7 win survives the family shift; only the precision-recall trade-off shifts.

\textbf{Decision.} KEEP as a robustness result. The cascade configuration is unchanged.

\subsection{Cross-cutting findings from the G-series}

Three patterns emerged across the eight phases that are independently worth noting:

\paragraph{(F1) The cascade is COMPOSITION-fragile, not COMPONENT-fragile.}
G3 is the cleanest demonstration: a +404\% rel standalone Hit@1 win broke cascade integration through two specific mechanisms (overlap-rerank vote shifts; L4 stacker score-scale mismatch). The structural assumption the cascade currently makes is that all retrievers produce comparable score distributions; violating this breaks the L3 free-novelty trigger and shifts the overlap-rerank vote.

\paragraph{(F2) Three independent rerankers all failed.}
G2 (cross-encoder reranker), Phase E (agent re-rank, pre-G-series), and G5 (LLM judge reranker) are three orthogonal designs with three negative results in the same direction. The cascade's overlap-rerank is empirically dominant for top-1 selection on this dataset; multi-retriever consensus beats single-LLM-judgment.

\paragraph{(F3) The retrieval channel was already at its local ceiling.}
Hit@5 stayed at $0.9124$ across the entire G-series. The actual win is in the L3 novelty channel (\texttt{novel\_recall} $+388\%$ rel). The original project framing as ``a retrieval improvement project'' should be reframed as ``the L3 novelty channel was underexploited.''

These three findings reshape the paper's framing. The detailed cross-phase synthesis is preserved in \texttt{docs4/META-ANALYSIS.md}.
